% ============================================================
% Week 3 — 16 Mar – 20 Mar 2026
% Compile:  cd report && make week_03
% ============================================================
\documentclass[11pt, a4paper]{article}

\usepackage[margin=2.5cm]{geometry}
\usepackage{graphicx}
\usepackage{booktabs}
\usepackage{multirow}
\usepackage{array}
\usepackage{xcolor}
\usepackage{hyperref}
\usepackage{amsmath}
\usepackage{amssymb}
\usepackage{caption}
\usepackage{subcaption}
\usepackage{float}
\usepackage{enumitem}
\usepackage{fancyhdr}
\usepackage{titlesec}
\usepackage{parskip}

\definecolor{sota}{HTML}{2E7D32}
\definecolor{good}{HTML}{1565C0}
\definecolor{warn}{HTML}{E65100}

\setlength{\headheight}{14pt}
\pagestyle{fancy}
\fancyhf{}
\lhead{\textbf{ESS Research} --- Week 3 Update}
\rhead{Inference-Time Prompt Gap in SAM3 Fine-tuning}
\cfoot{\thepage}

\titleformat{\section}{\large\bfseries}{}{0em}{}[\titlerule]
\titleformat{\subsection}{\normalsize\bfseries}{}{0em}{}

\hypersetup{colorlinks=true, linkcolor=blue, urlcolor=blue}
\newcolumntype{C}[1]{>{\centering\arraybackslash}p{#1}}

\begin{document}

% ── Cover ────────────────────────────────────────────────────
\begin{center}
    {\LARGE \textbf{Closing the Inference-Time Prompt Gap}}\\[0.3em]
    {\large in SAM3-Based Surgical Instrument Segmentation via Temporal Propagation}\\[0.5em]
    \rule{0.5\linewidth}{0.4pt}\\[0.4em]
    {\normalsize Weekly Research Progress Report --- \textbf{Week 3}}\\[0.3em]
    {\normalsize Week dates: 16 Mar -- 20 Mar 2026 \quad|\quad Report date: 20 March 2026}\\[0.3em]
    {\normalsize \href{https://github.com/skhan61/ess-research}{\texttt{github.com/skhan61/ess-research}}}
\end{center}
\vspace{0.5em}

\noindent Week~1 showed that SAM3+LoRA+box exceeds task-specific SOTA, and that
spatial prompt quality is the primary performance bottleneck.
Week~2 proposed robot kinematics as a GT source --- a direction requiring lab
dataset access that is not yet available.
Week~3 \textbf{pivots to a self-contained publishable direction} using only the
public UW-Sinus-Surgery-CL dataset already in hand:
identifying and closing the \textbf{inference-time prompt gap} in SAM3 fine-tuning.

\hrule
\vspace{1em}

% ============================================================
\section{Research Direction: What and Why}
% ============================================================

\subsection{Supervisor Guidance (Prof.\ Lee, Feb 25 2026)}

\begin{quote}
\textit{``I think your plan is very doable. But my suggestion is tune down the tone
to ONE paper. That means you will identify a specific algorithm, and improve from
there. You can start to use public dataset, and we have both cadaver and real
clinical data that contains high accuracy ($<$1\,mm error) ground truth.''}
\end{quote}

This week operationalises that guidance:
\begin{itemize}[noitemsep]
    \item \textbf{One algorithm:} SAM3+LoRA (already running).
    \item \textbf{Identify the failure:} the best prompt (GT box) is circular ---
          it requires the GT mask you are trying to predict. It cannot be
          reproduced at deployment.
    \item \textbf{Improve from there:} replace per-frame GT prompts with
          prompt-once temporal propagation using SAM3's built-in video memory.
    \item \textbf{Public dataset:} UW-Sinus-Surgery-CL. Lab data used later
          as extension or validation.
\end{itemize}

\subsection{The Circular Dependency}

In the H4 experiments (Week~1), the box prompt was generated as follows:

\begin{enumerate}[noitemsep]
    \item Load the GT binary mask for each frame.
    \item Find the bounding box of all instrument pixels: $[x_1, y_1, x_2, y_2]$.
    \item Pass that box to SAM3 as the spatial prompt.
\end{enumerate}

This is circular: to generate the box you need the GT mask; the GT mask is
what you are trying to predict. The high Dice scores reported for box-prompted
experiments ($0.9593$ live fold1) are therefore \textbf{upper bounds under an
idealised condition that does not exist at deployment}.

\subsection{The Gap Already Measured in Week 1}

The gap between realistic and idealised performance is already in the data:

\begin{center}
\begin{tabular}{lcccc}
\toprule
\textbf{Prompt} & \textbf{C$\to$C} & \textbf{L$\to$L f1} &
    \textbf{L$\to$L f2} & \textbf{L$\to$L f3} \\
\midrule
Text (realistic baseline, H3) & 0.9329 & 0.9049 & 0.8389 & 0.8886 \\
GT box (upper bound, H4)      & 0.9476 & 0.9593 & 0.9320 & 0.9548 \\
\midrule
\textbf{Gap} ($\Delta$Dice)   & 0.015  & 0.054  & 0.093  & 0.066 \\
\bottomrule
\end{tabular}
\end{center}

\noindent The gap is as large as \textbf{0.093 Dice} (live fold2, the hardest session).
No new experiments are required to establish this gap --- it exists in results
already collected. The research question for Week~3 onward is:

\begin{center}
\textit{Can temporal propagation close this gap without requiring GT at inference?}
\end{center}

% ============================================================
\section{What the Literature Has Done}
% ============================================================

The inference-time prompt gap is a known problem. Prior approaches all
require an \textbf{additional model or component} to generate prompts automatically:

\begin{center}
\begin{tabular}{p{3.5cm} p{4cm} p{5cm}}
\toprule
\textbf{Work} & \textbf{Approach} & \textbf{Limitation} \\
\midrule
Wu et al.~\cite{wu2024pointtracking} (2024) &
    Separate online point tracker generates prompts from tracked sparse points &
    Requires an additional tracker; manual initialisation still needed \\
Surgical-DeSAM (2024) &
    DETR detector auto-generates bounding boxes &
    Depends on detector accuracy; extra model required \\
SurgicalSAM~\cite{yue2024surgicalsam} (2024) &
    Class prototype embeddings (no spatial prompt) &
    No spatial localisation; treats whole instrument as one entity \\
SAM2S (2025) &
    Enhanced memory + semantic long-term tracking &
    Requires large surgical video dataset (61k frames) for training \\
\bottomrule
\end{tabular}
\end{center}

\noindent \textbf{The gap:} no prior work uses SAM3's \textit{native} video memory
to propagate a single initial prompt across a clip --- without any additional
model, without per-frame GT, and with LoRA fine-tuning on a small domain-specific
dataset.

% ============================================================
\section{Proposed Fix: Prompt-Once Temporal Propagation}
% ============================================================

\subsection{Core Idea}

SAM3 inherits SAM2's video memory module~\cite{ravi2024sam2}. After seeing a
segmented frame, SAM3 stores a memory embedding of the object it found.
On subsequent frames it uses that memory to locate and segment the same object ---
without any new prompt.

The proposed inference protocol:

\begin{center}
\begin{tabular}{lp{9cm}}
\toprule
\textbf{Frame} & \textbf{Action} \\
\midrule
Frame 1 & GT box prompt $\to$ SAM3 segments instrument $\to$ stores in memory \\
Frames $2 \ldots N$ & No prompt. SAM3 reads memory, tracks instrument, predicts mask. \\
\bottomrule
\end{tabular}
\end{center}

GT is touched \textbf{once per video clip}, not once per frame.
In practice, ``prompting frame 1'' is feasible in a clinical workflow --- a
clinician or a simple click-interface identifies the instrument in the first frame.

\subsection{Why This Is Better Than Prior Work}

\begin{itemize}[noitemsep]
    \item No separate detector or tracker.
    \item No large-scale pre-training dataset required.
    \item SAM3's memory mechanism is already trained for exactly this use case.
    \item Directly comparable to Week~1 results: same model, same dataset,
          same fine-tuning --- only the prompting strategy changes.
\end{itemize}

\subsection{The Paper in One Sentence}

\begin{quote}
\textit{SAM3+LoRA fine-tuned with per-frame GT box prompts achieves SOTA on
UW-Sinus-Surgery-CL; we show this result is unreproducible at deployment due
to a circular GT dependency, quantify the resulting performance gap, and close
it via prompt-once temporal propagation using SAM3's native video memory.}
\end{quote}

% ============================================================
\section{Planned Experiments --- Week 3}
% ============================================================

Three experiments are needed. All use the existing SAM3+LoRA checkpoints
from Week~1 --- no retraining required.

\begin{center}
\begin{tabular}{C{0.8cm} p{3.5cm} p{5cm} p{3.5cm}}
\toprule
\textbf{ID} & \textbf{Name} & \textbf{What It Measures} & \textbf{Expected Result} \\
\midrule
H6 & Prompt-gap quantification &
    Re-run H3 (text) and H4 (GT box) on \textit{video sequences} (ordered frames),
    not shuffled frames. Establishes the sequence-level gap. &
    Gap $\sim$0.05--0.09 Dice, consistent with Week~1 \\
H7 & Temporal propagation &
    Prompt frame 1 of each clip with GT box. Frames 2--$N$: no prompt.
    Use SAM3 memory module. Evaluate Dice on all frames. &
    Dice between text (H3) and GT box (H4) upper bound \\
H8 & Propagation robustness &
    Vary clip length: 5, 10, 20, 50 frames. Measure Dice decay as clip grows.
    Identifies how many frames the memory stays reliable. &
    Graceful degradation; plateau or slow decay \\
\bottomrule
\end{tabular}
\end{center}

\subsection{Key Comparison Table (Target)}

After Week~3, the central result table will be:

\begin{center}
\begin{tabular}{lcccc}
\toprule
\textbf{Setting} & \textbf{C$\to$C} & \textbf{L$\to$L f1} &
    \textbf{L$\to$L f2} & \textbf{L$\to$L f3} \\
\midrule
Text prompt, per-frame (H3) --- \textit{realistic, weak} & 0.9329 & 0.9049 & 0.8389 & 0.8886 \\
GT box, per-frame (H4) --- \textit{upper bound, circular} & 0.9476 & 0.9593 & 0.9320 & 0.9548 \\
\textbf{Temporal propagation (H7)} --- \textit{realistic, strong} & \textbf{?} & \textbf{?} & \textbf{?} & \textbf{?} \\
\midrule
HFRF-Net (task-specific SOTA) & --- & 0.9374 & 0.9374 & 0.9374 \\
\bottomrule
\end{tabular}
\end{center}

\noindent If H7 lands above H3 and close to H4 across most folds, the contribution
is demonstrated: temporal propagation closes the prompt gap without requiring GT
per frame.

% ============================================================
\section{On the Broader Interpretation of Lee's Direction}
% ============================================================

Lee's instruction --- \textit{``identify a specific algorithm, improve from
there''} --- is interpreted as follows and is not misread:

\begin{enumerate}[noitemsep]
    \item \textbf{Specific algorithm:} SAM3+LoRA. Already chosen and running.
    \item \textbf{Identify the failure:} the inference-time prompt gap. Already
          quantified in Week~1 data.
    \item \textbf{Improve from there:} temporal propagation. Uses the same model
          with a different inference strategy.
    \item \textbf{Public dataset first:} UW-Sinus-Surgery-CL throughout.
          Lab data (cadaver + clinical, $<$1\,mm GT) is a natural extension
          once the core method is validated.
    \item \textbf{Generality:} the prompt gap problem exists for all SAM-based
          methods that report GT-box evaluation (SurgSAM-2, SurgicalSAM,
          wu2024pointtracking). The fix applies beyond SAM3 --- any model
          with a video memory module can use the same protocol.
\end{enumerate}

% ============================================================
\section{Repository}
% ============================================================

\begin{center}
\begin{tabular}{ll}
\toprule
\textbf{Item} & \textbf{Detail} \\
\midrule
Repository   & \texttt{github.com/skhan61/ess-research} \\
Entry point  & \texttt{uv run python scripts/run\_all\_experiments.py --plan research\_plan.yaml} \\
Debug flag   & \texttt{--debug} \\
\bottomrule
\end{tabular}
\end{center}

\nocite{qin2020uw,hfrfnet2025,ravi2024sam2,yue2024surgicalsam,wu2024pointtracking,liu2024surgsam2}
\bibliographystyle{plain}
\bibliography{../references}

\end{document}
